% Options for packages loaded elsewhere
\PassOptionsToPackage{unicode}{hyperref}
\PassOptionsToPackage{hyphens}{url}
\documentclass[
  11pt,
]{memoir}
\usepackage{xcolor}
\usepackage[margin=1in]{geometry}
\usepackage{amsmath,amssymb}
\setcounter{secnumdepth}{5}
\usepackage{iftex}
\ifPDFTeX
  \usepackage[T1]{fontenc}
  \usepackage[utf8]{inputenc}
  \usepackage{textcomp} % provide euro and other symbols
\else % if luatex or xetex
  \usepackage{unicode-math} % this also loads fontspec
  \defaultfontfeatures{Scale=MatchLowercase}
  \defaultfontfeatures[\rmfamily]{Ligatures=TeX,Scale=1}
\fi
\usepackage{lmodern}
\ifPDFTeX\else
  % xetex/luatex font selection
\fi
% Use upquote if available, for straight quotes in verbatim environments
\IfFileExists{upquote.sty}{\usepackage{upquote}}{}
\IfFileExists{microtype.sty}{% use microtype if available
  \usepackage[]{microtype}
  \UseMicrotypeSet[protrusion]{basicmath} % disable protrusion for tt fonts
}{}
\makeatletter
\@ifundefined{KOMAClassName}{% if non-KOMA class
  \IfFileExists{parskip.sty}{%
    \usepackage{parskip}
  }{% else
    \setlength{\parindent}{0pt}
    \setlength{\parskip}{6pt plus 2pt minus 1pt}}
}{% if KOMA class
  \KOMAoptions{parskip=half}}
\makeatother
\usepackage{color}
\usepackage{fancyvrb}
\newcommand{\VerbBar}{|}
\newcommand{\VERB}{\Verb[commandchars=\\\{\}]}
\DefineVerbatimEnvironment{Highlighting}{Verbatim}{commandchars=\\\{\}}
% Add ',fontsize=\small' for more characters per line
\newenvironment{Shaded}{}{}
\newcommand{\AlertTok}[1]{\textcolor[rgb]{1.00,0.00,0.00}{\textbf{#1}}}
\newcommand{\AnnotationTok}[1]{\textcolor[rgb]{0.38,0.63,0.69}{\textbf{\textit{#1}}}}
\newcommand{\AttributeTok}[1]{\textcolor[rgb]{0.49,0.56,0.16}{#1}}
\newcommand{\BaseNTok}[1]{\textcolor[rgb]{0.25,0.63,0.44}{#1}}
\newcommand{\BuiltInTok}[1]{\textcolor[rgb]{0.00,0.50,0.00}{#1}}
\newcommand{\CharTok}[1]{\textcolor[rgb]{0.25,0.44,0.63}{#1}}
\newcommand{\CommentTok}[1]{\textcolor[rgb]{0.38,0.63,0.69}{\textit{#1}}}
\newcommand{\CommentVarTok}[1]{\textcolor[rgb]{0.38,0.63,0.69}{\textbf{\textit{#1}}}}
\newcommand{\ConstantTok}[1]{\textcolor[rgb]{0.53,0.00,0.00}{#1}}
\newcommand{\ControlFlowTok}[1]{\textcolor[rgb]{0.00,0.44,0.13}{\textbf{#1}}}
\newcommand{\DataTypeTok}[1]{\textcolor[rgb]{0.56,0.13,0.00}{#1}}
\newcommand{\DecValTok}[1]{\textcolor[rgb]{0.25,0.63,0.44}{#1}}
\newcommand{\DocumentationTok}[1]{\textcolor[rgb]{0.73,0.13,0.13}{\textit{#1}}}
\newcommand{\ErrorTok}[1]{\textcolor[rgb]{1.00,0.00,0.00}{\textbf{#1}}}
\newcommand{\ExtensionTok}[1]{#1}
\newcommand{\FloatTok}[1]{\textcolor[rgb]{0.25,0.63,0.44}{#1}}
\newcommand{\FunctionTok}[1]{\textcolor[rgb]{0.02,0.16,0.49}{#1}}
\newcommand{\ImportTok}[1]{\textcolor[rgb]{0.00,0.50,0.00}{\textbf{#1}}}
\newcommand{\InformationTok}[1]{\textcolor[rgb]{0.38,0.63,0.69}{\textbf{\textit{#1}}}}
\newcommand{\KeywordTok}[1]{\textcolor[rgb]{0.00,0.44,0.13}{\textbf{#1}}}
\newcommand{\NormalTok}[1]{#1}
\newcommand{\OperatorTok}[1]{\textcolor[rgb]{0.40,0.40,0.40}{#1}}
\newcommand{\OtherTok}[1]{\textcolor[rgb]{0.00,0.44,0.13}{#1}}
\newcommand{\PreprocessorTok}[1]{\textcolor[rgb]{0.74,0.48,0.00}{#1}}
\newcommand{\RegionMarkerTok}[1]{#1}
\newcommand{\SpecialCharTok}[1]{\textcolor[rgb]{0.25,0.44,0.63}{#1}}
\newcommand{\SpecialStringTok}[1]{\textcolor[rgb]{0.73,0.40,0.53}{#1}}
\newcommand{\StringTok}[1]{\textcolor[rgb]{0.25,0.44,0.63}{#1}}
\newcommand{\VariableTok}[1]{\textcolor[rgb]{0.10,0.09,0.49}{#1}}
\newcommand{\VerbatimStringTok}[1]{\textcolor[rgb]{0.25,0.44,0.63}{#1}}
\newcommand{\WarningTok}[1]{\textcolor[rgb]{0.38,0.63,0.69}{\textbf{\textit{#1}}}}
\usepackage{longtable,booktabs,array}
\usepackage{calc}
\newcommand{\real}[1]{#1} % for calculating minipage widths
% Correct order of tables after \paragraph or \subparagraph
\usepackage{etoolbox}
\makeatletter
\patchcmd\longtable{\par}{\if@noskipsec\mbox{}\fi\par}{}{}
\makeatother
% Allow footnotes in longtable head/foot
\IfFileExists{footnotehyper.sty}{\usepackage{footnotehyper}}{\usepackage{footnote}}
\makesavenoteenv{longtable}
\usepackage{graphicx}
\makeatletter
\newsavebox\pandoc@box
\newcommand*\pandocbounded[1]{% scales image to fit in text height/width
  \sbox\pandoc@box{#1}%
  \Gscale@div\@tempa{\textheight}{\dimexpr\ht\pandoc@box+\dp\pandoc@box\relax}%
  \Gscale@div\@tempb{\linewidth}{\wd\pandoc@box}%
  \ifdim\@tempb\p@<\@tempa\p@\let\@tempa\@tempb\fi% select the smaller of both
  \ifdim\@tempa\p@<\p@\scalebox{\@tempa}{\usebox\pandoc@box}%
  \else\usebox{\pandoc@box}%
  \fi%
}
% Set default figure placement to htbp
\def\fps@figure{htbp}
\makeatother
\setlength{\emergencystretch}{3em} % prevent overfull lines
\providecommand{\tightlist}{%
  \setlength{\itemsep}{0pt}\setlength{\parskip}{0pt}}
% Memoir PDF tweaks for REPORT.md script tables and figures.
\usepackage{etoolbox}
\usepackage{xurl}
\usepackage{graphicx}

\AtBeginEnvironment{longtable}{%
  \small
  \setlength{\tabcolsep}{4pt}%
  \sloppy
}
\usepackage{bookmark}
\IfFileExists{xurl.sty}{\usepackage{xurl}}{} % add URL line breaks if available
\urlstyle{same}
\hypersetup{
  hidelinks,
  pdfcreator={LaTeX via pandoc}}

\author{}
\date{}

\begin{document}
\frontmatter

\mainmatter
\chapter{Identification, Generalization, and Dense Associative Memory on
Grid Codes and Naturalistic
Vision}\label{identification-generalization-and-dense-associative-memory-on-grid-codes-and-naturalistic-vision}

\textbf{Ahmed Al Sunbati}\\
Neukom Scholar project handoff report\\
Dartmouth College, June 2026

\begin{center}\rule{0.5\linewidth}{0.5pt}\end{center}

\section{Introduction}\label{introduction}

The starting question for this term was whether we could replicate a
meaninguflly tradeoff between \textbf{identification} and
\textbf{generalization} in vision systems coupled with DAM.
Identification means: store a pattern, cue with a noisy version of the
same pattern, and ask whether retrieval lands on the correct stored
item. Generalization means: cue with a novel pattern and ask whether
retrieval lands on the stored item that best matches it under some
similarity rule.

I built two parallel tracks. The \textbf{grid track}, which was finished
last term and optimized over the start of the spring to reduce runtime
to under 5 mins per experiment, uses a synthetic 2D spatial code
(inspired by grid cells) where the ident/gen tradeoff is real and easy
to measure. The \textbf{vision track} uses pretrained image models (ViT,
CLIP, ConvNeXt) on colored squares and on real Peterson images with
human similarity ratings. Everything needed to rerun or reload results
lives in the \texttt{handoff/} folder: two notebooks, CLI scripts,
bundled data, and precomputed outputs under \texttt{results/}.

\begin{center}\rule{0.5\linewidth}{0.5pt}\end{center}

\section{1. Core memory model: Dense Associative
Memory}\label{core-memory-model-dense-associative-memory}

Both tracks share the same memory module, implemented in
\texttt{core/DAM.py} as \texttt{DenseAssociativeMemory}.

\subsection{1.1 Identification and
generalization}\label{identification-and-generalization}

\begin{quote}
\textbf{Identification:} retrieve the stored pattern from a noisy
encoding of the same point.\\
\textbf{Generalization:} retrieve the nearest stored pattern from a
novel test point.
\end{quote}

On the grid, ``novel'' means a fresh draw from the 2D input space. On
vision, identification probes are corrupted copies of stored images;
generalization probes are held-out images whose correct target is the
stored image with highest human similarity (naturalistic) or highest
task similarity (synthetic).

\subsection{1.2 Energy and retrieval}\label{energy-and-retrieval}

Classical Hopfield networks use a quadratic energy in the overlap
between state and stored patterns. Dense Associative Memory raises that
overlap to order $n > 2$ inside a rectified polynomial,
which increases effective capacity. The energy for state vector
$\mathbf{s}$ and stored patterns $\mathbf{p}_{\mu}$ is:

\[
E(\mathbf{s}) = -\sum_{\mu} F(\mathbf{p}_\mu \cdot \mathbf{s}) + \frac{\lambda}{2}\|\mathbf{s}\|^2, \qquad F(x) = [\max(0, x)]^n
\]

Retrieval starts from a noisy cue and runs asynchronous neuron updates.
Each step moves one coordinate of $\mathbf{s}$ toward a value that
lowers the energy, using inverse temperature $\beta$, relaxation rate
$\alpha$, and optional L2 regularization $\lambda$. The network settles
on whichever stored pattern has the strongest overlap with the final
state.

For large sweeps I used the \textbf{Numba} backend
(\texttt{retrieve\_differential}) because the pure NumPy loop is too
slow when you are crossing thousands of ($(c, n, K)$) combinations.

\begin{center}\rule{0.5\linewidth}{0.5pt}\end{center}

\section{2. Grid encoder and synthetic spatial
benchmark}\label{grid-encoder-and-synthetic-spatial-benchmark}

The grid track maps 2D coordinates into high-dimensional patterns before
they enter DAM. \texttt{core/Encoder.py} implements
\texttt{GridEncoder}, modeled after entorhinal grid cells.

\subsection{2.1 Encoding}\label{encoding}

Each module (m) runs at spatial frequency proportional to
\texttt{base\_freq\ ·\ c\^{}m}, where $c$ is the \textbf{scale factor}
swept in the multiscale experiment. Within a module there are several
orientation vectors $\mathbf{k}_{o}$ in the plane and several
phase-shifted grid cells. A 2D position $\mathbf{x}$ is projected onto
each grid:

\[
\phi_{m,o,p}(\mathbf{x}) = 2\pi f_m (\mathbf{x} \cdot \mathbf{k}_o) + \psi_{m,o,p}
\]

The encoded vector stacks $\cos\phi$ and $\sin\phi$ over all modules,
orientations, and cells. Nearby locations get similar codes; far apart
locations diverge. That gives DAM a smooth spatial memory substrate
where we control geometry directly.

Default encoder settings in the notebook (4 modules, 3 orientations, 4
cells per orientation, \texttt{base\_freq\ =\ 0.5}) match
\texttt{lib/grid\_defaults.py}.

\begin{center}\rule{0.5\linewidth}{0.5pt}\end{center}

\section{\texorpdfstring{3. Notebook:
\texttt{GridDAM.ipynb}}{3. Notebook: GridDAM.ipynb}}\label{notebook-griddam.ipynb}

\subsection{3.1 Setup}\label{setup}

The setup cell bootstraps paths from \texttt{handoff/lib/paths.py} and
sets:

\begin{longtable}[]{@{}
  >{\raggedright\arraybackslash}p{(\linewidth - 4\tabcolsep) * \real{0.2500}}
  >{\raggedright\arraybackslash}p{(\linewidth - 4\tabcolsep) * \real{0.3750}}
  >{\raggedright\arraybackslash}p{(\linewidth - 4\tabcolsep) * \real{0.3750}}@{}}
\toprule\noalign{}
\begin{minipage}[b]{\linewidth}\raggedright
Flag
\end{minipage} & \begin{minipage}[b]{\linewidth}\raggedright
Default
\end{minipage} & \begin{minipage}[b]{\linewidth}\raggedright
Meaning
\end{minipage} \\
\midrule\noalign{}
\endhead
\bottomrule\noalign{}
\endlastfoot
\texttt{QUICK\_MODE} & \texttt{False} & Smaller parameter grids for
smoke tests \\
\texttt{LOAD\_EXISTING} & \texttt{False} & Skip multiscale sweep if
{\ttfamily\nolinkurl{results/grid/multiscale/ident\_results.json}} already exists \\
\texttt{OUTPUT\_DIR} & {\ttfamily\nolinkurl{results/grid}} & All grid outputs \\
\end{longtable}

\subsection{3.2 Structure preservation}\label{structure-preservation}

Before any DAM sweep, the notebook checks that \textbf{encoded}
distances still correlate with \textbf{2D} distances as $c$ changes. If
the code stops preserving geometry, ident/gen numbers are hard to
interpret.

\begin{quote}
Before running DAM sweeps, check that encoded distances still correlate
with 2D distances as \texttt{c} changes.
\end{quote}

\subsection{3.3 Multiscale ident/gen
sweep}\label{multiscale-identgen-sweep}

The pipeline, quoted from the notebook:

\begin{quote}
\begin{enumerate}
\def\labelenumi{\arabic{enumi}.}
\tightlist
\item
  \textbf{Encoder:} build a multiscale grid code with
  \texttt{n\_modules} modules. Module \texttt{m} uses frequency
  \texttt{base\_freq\ *\ c\^{}m}, where \texttt{c} is the scale
  factor.\\
\item
  \textbf{Memory:} store \texttt{K} encoded 2D points in a Dense
  Associative Memory. Vary the energy order \texttt{n} of the DAM.\\
\item
  \textbf{Identification trials:} cue with noisy encodings of stored
  points.\\
\item
  \textbf{Generalization trials:} cue with encodings of novel test
  points.\\
\item
  \textbf{Sweep} over scale factor \texttt{c}, DAM order \texttt{n}, and
  number stored \texttt{K}. Plot identification vs generalization.
\end{enumerate}
\end{quote}

\textbf{Main finding:} generalization peaks near scale factor $c \approx \sqrt{e} \approx 1.65$ while identification keeps improving at
larger $c$.

\begin{figure}
\centering
\pandocbounded{\includegraphics[keepaspectratio,alt={Figure 1. Multiscale tradeoff curves. Identification and generalization accuracy vs scale factor c for several DAM orders n.}]{results/grid/multiscale/tradeoff_by_n.png}}
\caption{Figure 1. Multiscale tradeoff curves. Identification and
generalization accuracy vs scale factor c for several DAM orders n.}
\end{figure}

Outputs: \texttt{ident\_results.json}, \texttt{gen\_results.json},
\texttt{config.json}, and the tradeoff figure above.

\subsection{3.4 Encoder sensitivity
sweeps}\label{encoder-sensitivity-sweeps}

This part varies \textbf{one encoder parameter at a time} (orientations,
base frequency, module count, cells per orientation, or 2D sampling
distribution) and tracks the gen-optimal scale factor (c\^{}*) at fixed
DAM order (n = 4).

\begin{quote}
Change one encoder parameter at a time and track the gen-optimal scale
factor \texttt{c} at \texttt{n=4}.
\end{quote}

Each experiment writes JSON summaries and tradeoff PNGs under
{\ttfamily\nolinkurl{results/grid/breakit/\{experiment\}}/}.

\begin{figure}
\centering
\pandocbounded{\includegraphics[keepaspectratio,alt={Figure 2. Example break-it tradeoff (orientations = 6).}]{results/grid/breakit/orientations/orientations_6_tradeoff.png}}
\caption{Figure 2. Example break-it tradeoff (orientations = 6).}
\end{figure}

\begin{center}\rule{0.5\linewidth}{0.5pt}\end{center}

\section{\texorpdfstring{4. Notebook:
\texttt{VisionExperiments.ipynb}}{4. Notebook: VisionExperiments.ipynb}}\label{notebook-visionexperiments.ipynb}

\subsection{4.1 Pipeline overview}\label{pipeline-overview}

\begin{quote}
\begin{enumerate}
\def\labelenumi{\arabic{enumi}.}
\tightlist
\item
  Synthetic baseline on colored squares\\
\item
  Multi-model comparison on the synthetic task\\
\item
  Naturalistic baseline on real images with human similarity\\
\item
  Dense Associative Memory on the naturalistic benchmark\\
\item
  ViT vs CLIP on the Animals dataset: \textbf{layer-index} and
  \textbf{\S{}distance-rank}
\end{enumerate}
\end{quote}

\textbf{Notebook flags:}

\begin{longtable}[]{@{}
  >{\raggedright\arraybackslash}p{(\linewidth - 4\tabcolsep) * \real{0.2500}}
  >{\raggedright\arraybackslash}p{(\linewidth - 4\tabcolsep) * \real{0.3750}}
  >{\raggedright\arraybackslash}p{(\linewidth - 4\tabcolsep) * \real{0.3750}}@{}}
\toprule\noalign{}
\begin{minipage}[b]{\linewidth}\raggedright
Flag
\end{minipage} & \begin{minipage}[b]{\linewidth}\raggedright
Default
\end{minipage} & \begin{minipage}[b]{\linewidth}\raggedright
Meaning
\end{minipage} \\
\midrule\noalign{}
\endhead
\bottomrule\noalign{}
\endlastfoot
\texttt{RUN\_EXPERIMENTS} & \texttt{False} & Reload bundled JSON/PNG
from \texttt{results/} \\
\texttt{RUN\_EXPERIMENTS} & \texttt{True} & Regenerate into the same
paths (slow on CPU) \\
\texttt{QUICK\_MODE} & use with live runs & Smaller grids: synthetic
caps; DAM quick = animals only, 1 stage-A config \\
\texttt{QUICK\_MODE=False} (DAM live) & --- & All three categories, 6
stage-A configs, stage B skipped, 1 seed, storage 40, ViT+CLIP \\
\end{longtable}

\subsection{4.2 Identification corruption
protocol}\label{identification-corruption-protocol}

Tracks 1--4 (synthetic baseline, model comparison, naturalistic
baseline, naturalistic DAM) share a \textbf{stacked} identification
protocol:

\begin{longtable}[]{@{}ll@{}}
\toprule\noalign{}
Component & Setting \\
\midrule\noalign{}
\endhead
\bottomrule\noalign{}
\endlastfoot
Pixels & noise (std = 5) + shift (\(\leq\) 2 px) + 50\% occlusion \\
Retrieval scores & decision noise std = 0.01 on ident trials only \\
\end{longtable}

\subsection{4.3 Stage 1: Synthetic
baseline}\label{stage-1-synthetic-baseline}

Colored squares at random positions; layerwise ViT features; cosine
retrieval. Produces a line plot of identification vs generalization
accuracy across transformer layers. Early debugging found that z-score
normalization on tiny synthetic batches broke ident sanity checks;
switching to \textbf{L2-only} normalization fixed that.

\subsection{4.4 Stage 2: Model comparison (color vs
location)}\label{stage-2-model-comparison-color-vs-location}

What if only color or only location varies? Three task modes are swept
per backbone:

\begin{itemize}
\tightlist
\item
  \texttt{mixed\_color\_position} (both vary)
\item
  \texttt{color\_only} (square fixed position, color jitter)
\item
  \texttt{position\_only} (fixed color, position varies)
\end{itemize}

Backbones in the bundled summary: ViT, CLIP ViT, ConvNeXt. The notebook
plots best-generalization, best-identification, and best-margin configs
per model in ident/gen space.

\begin{quote}
Synthetic model-comparison summary from
{\ttfamily\nolinkurl{results/vision/model\_comparison/}}. Each backbone contributes
three highlighted configs: best generalization, best identification, and
best margin.
\end{quote}

\subsection{4.5 Stage 3: Naturalistic
baseline}\label{stage-3-naturalistic-baseline}

Peterson dataset: real photos plus a human similarity matrix per
category (fruits, vegetables, animals). Generalization accuracy =
percent of probes where retrieval picks the stored image with
\textbf{highest human-rated similarity}. Vegetables uses
\texttt{required\_exemplars\ =\ 3} internally so leave-one-exemplar-out
folds stay balanced.

\begin{quote}
Real-image retrieval with human similarity for generalization.
Identification probes use stacked \texttt{noise\_shift\_occlusion} plus
decision noise.
\end{quote}

\subsection{4.6 Stage 4: Naturalistic DAM (ViT vs
CLIP)}\label{stage-4-naturalistic-dam-vit-vs-clip}

Same task definitions as baseline, but features pass through DAM with a
small config search.

The notebook bar chart reads \texttt{vit\_vs\_clip\_head\_to\_head.json}
and plots \textbf{gen accuracy delta} (DAM minus baseline) for matched
ViT vs CLIP branches. There are \textbf{seven} branch pairs (category ×
pooling × \texttt{anchor\_kind}, where \texttt{anchor\_kind} is
\texttt{retrieval} or \texttt{alignment}), so labels like
\texttt{animals-mean\_tokens} can appear twice when retrieval and
alignment anchors differ.

\begin{quote}
Gen accuracy delta = DAM minus baseline on the human-similarity
generalization task (\% of probes where retrieval picks the human-rated
best stored match; positive = DAM helps).
\end{quote}

\subsection{4.7 Stage 5: ViT vs CLIP on Animals
dataset}\label{stage-5-vit-vs-clip-on-animals-dataset}

Two views of the same metrics on the animals category only:

\textbf{layer-index:} points connected in transformer layer order;
labels \texttt{L0}\ldots{}\texttt{L11}. Reload from
{\ttfamily\nolinkurl{results/image\_transformers/\{setting\}}/}.

\textbf{distance-rank:} same axes, but points ordered by ascending mean
pairwise cosine distance among stored features; labels
\texttt{R0}\ldots{}\texttt{R11}. Reload from
{\ttfamily\nolinkurl{results/image\_transformers\_similarity\_ordered\_fixed/\{setting\}}/}.

Six settings: \texttt{easy\_s40}, \texttt{easy\_s80},
\texttt{occ50\_s40}, \texttt{occ50\_s80},
\texttt{occ50\_dnoise001\_s40}, \texttt{occ50\_dnoise001\_s80}. Each
folder has five PNGs: baseline plus DAM orders (n = 2, 4, 6, 8).

\begin{figure}
\centering
\pandocbounded{\includegraphics[keepaspectratio,alt={Figure 3. occ50\_s40 baseline, layer-index order.}]{results/image_transformers/occ50_s40/baseline_layers_gen_vs_ident.png}}
\caption{Figure 3. occ50\_s40 baseline, layer-index order.}
\end{figure}

\begin{figure}
\centering
\pandocbounded{\includegraphics[keepaspectratio,alt={Figure 4. occ50\_s40 baseline, distance-rank order.}]{results/image_transformers_similarity_ordered_fixed/occ50_s40/baseline_layers_gen_vs_ident_ordered_by_distance.png}}
\caption{Figure 4. occ50\_s40 baseline, distance-rank order.}
\end{figure}

\begin{quote}
Same axes as \S{}5A, but point sequence follows \textbf{ascending mean
pairwise cosine distance} (stored-feature geometry).
\end{quote}

Thereordering is \textbf{visualization only}; it does not change the
underlying metrics (see
{\ttfamily\nolinkurl{results/image\_transformers\_similarity\_ordered\_fixed/overall\_insights.md}}).
The reordering is just a different view of the same layerwise
represenation geometry.

\begin{center}\rule{0.5\linewidth}{0.5pt}\end{center}

\section{5. Results summary}\label{results-summary}

\subsection{5.1 Grid (multiscale)}\label{grid-multiscale}

On the grid task the ident/gen tradeoff is clear: generalization peaks
near $c \approx 1.65$, identification rises with larger $c$. Break-it
sweeps show how encoder parameters shift that optimum.

\subsection{5.2 Vision: synthetic}\label{vision-synthetic}

Synthetic colored squares were useful for debugging.

\begin{itemize}
\tightlist
\item
  After the L2-only preprocessing fix, sanity checks passed.
\item
  No convincing ident/gen tradeoff appeared on synthetic features.
\item
  Color was easy; position was weaker.
\item
  Swapping backbone (ViT, DeiT, ConvNeXt, ResNet, CLIP) did not fix it.
\item
  Splitting \textbf{color-only} vs \textbf{position-only} tasks showed
  color dominating; position alone was especially weak.
\end{itemize}

So the synthetic distribution (colored squares at random locations) is
probably too unusual relative to ImageNet pretraining, which matches
your intution from email.

\subsection{5.3 Vision: naturalistic
baseline}\label{vision-naturalistic-baseline}

Moving to Peterson images changed everything. Across fruits, vegetables,
and animals:

\begin{itemize}
\tightlist
\item
  \textbf{Late layers} gave the best retrieval accuracy.
\item
  \textbf{Mid layers} gave the best alignment with human similarity
  (Spearman correlation between model RDM and human RDM).
\end{itemize}

Best baseline generalization (ViT, bundled notebook regen, stacked ident
protocol):

\begin{longtable}[]{@{}ll@{}}
\toprule\noalign{}
Category & Best ViT gen accuracy \\
\midrule\noalign{}
\endhead
\bottomrule\noalign{}
\endlastfoot
Fruits & 55.0\% \\
Vegetables & 53.8\% \\
Animals & 65.0\% \\
\end{longtable}

(CLIP image-tower retrieval peaks slightly higher on fruits at 57.5\%;
see 5.4.)

Best human-alignment Spearman (ViT mid-layer rows, not layer 11):

\begin{longtable}[]{@{}ll@{}}
\toprule\noalign{}
Category & human\_rdm\_spearman \\
\midrule\noalign{}
\endhead
\bottomrule\noalign{}
\endlastfoot
Fruits & 0.379 \\
Vegetables & 0.462 \\
Animals & 0.423 \\
\end{longtable}

Retrieval and alignment \textbf{do not peak at the same layer}. \#\#\#
5.4 CLIP as encoder (image only)

CLIP (ViT image tower, no text) boosted alignment a lot:

\begin{longtable}[]{@{}ll@{}}
\toprule\noalign{}
Category & Best CLIP alignment \\
\midrule\noalign{}
\endhead
\bottomrule\noalign{}
\endlastfoot
Fruits & 0.580 \\
Vegetables & 0.494 \\
Animals & 0.550 \\
\end{longtable}

But ViT still owned the strongest top retrieval rows. \textbf{ViT for
retrieval, CLIP for alignment.}

\subsection{5.5 DAM on naturalistic
features}\label{dam-on-naturalistic-features}

After baseline retrieval on real Peterson images, DAM was applied on top
of those features. This section summarizes whether that heled retrieval
in generalization and whether ViT or CLIP was the better backbone.

DAM helped unevenly. Numbers below come from the \textbf{capped notebook
regen} in {\ttfamily\nolinkurl{results/vision/naturalistic\_dam/}}. A separate full
CLI sweep can explore larger config grids.

\begin{longtable}[]{@{}
  >{\raggedright\arraybackslash}p{(\linewidth - 2\tabcolsep) * \real{0.2703}}
  >{\raggedright\arraybackslash}p{(\linewidth - 2\tabcolsep) * \real{0.7297}}@{}}
\toprule\noalign{}
\begin{minipage}[b]{\linewidth}\raggedright
Category
\end{minipage} & \begin{minipage}[b]{\linewidth}\raggedright
DAM story (bundled regen)
\end{minipage} \\
\midrule\noalign{}
\endhead
\bottomrule\noalign{}
\endlastfoot
Fruits & 0 qualifying wins; raw gen deltas up to +7.5 pp (CLIP
\texttt{mean\_tokens} retrieval) but probe-RDM alignment drops \\
Vegetables & 0 qualifying wins; raw deltas up to +7.7 pp (ViT
\texttt{mean\_tokens} retrieval) with modest alignment lift \\
Animals & 6 qualifying wins; best clean branch 52.5\% → 55.0\% gen (+2.5
pp, ViT \texttt{cls} retrieval, layer 11) \\
\end{longtable}

Head-to-head gen deltas (same capped run):

\begin{longtable}[]{@{}lll@{}}
\toprule\noalign{}
Branch & ViT Δgen & CLIP Δgen \\
\midrule\noalign{}
\endhead
\bottomrule\noalign{}
\endlastfoot
animals cls retrieval & +2.5 pp & +2.5 pp \\
animals mean\_tokens retrieval & +2.5 pp & +1.2 pp \\
animals mean\_tokens alignment & −1.2 pp & −1.3 pp \\
fruits mean\_tokens retrieval & +2.5 pp & \textbf{+7.5 pp} \\
vegetables mean\_tokens retrieval & \textbf{+7.7 pp} & +5.1 pp \\
\end{longtable}

So animals shows the most \textbf{stable} DAM wins under the qualifying
rule (improves generalization without breaking memory or geometry).

Matched ViT vs CLIP under DAM: \textbf{ViT wins on vegetables and
animals retrieval branches}; \textbf{CLIP wins on fruits
\texttt{mean\_tokens} retrieval}.

\subsection{5.6 Grid vs vision
(comparison)}\label{grid-vs-vision-comparison}

\begin{longtable}[]{@{}
  >{\raggedright\arraybackslash}p{(\linewidth - 4\tabcolsep) * \real{0.0714}}
  >{\raggedright\arraybackslash}p{(\linewidth - 4\tabcolsep) * \real{0.4286}}
  >{\raggedright\arraybackslash}p{(\linewidth - 4\tabcolsep) * \real{0.5000}}@{}}
\toprule\noalign{}
\begin{minipage}[b]{\linewidth}\raggedright
\end{minipage} & \begin{minipage}[b]{\linewidth}\raggedright
Grid track
\end{minipage} & \begin{minipage}[b]{\linewidth}\raggedright
Vision track
\end{minipage} \\
\midrule\noalign{}
\endhead
\bottomrule\noalign{}
\endlastfoot
Stimulus & Synthetic 2D points & Colored squares, then real photos \\
Features & GridEncoder & ViT / CLIP / ConvNeXt layers \\
Gen target & Nearest stored in 2D & Task similarity or human
similarity \\
Ident/gen tradeoff? & \textbf{Yes} (multiscale sweep) & \textbf{No} on
synthetic; different story on naturalistic \\
Main positive result & $c \approx \sqrt{e}$ gen peak & Retrieval vs
alignment split; DAM on animals \\
\end{longtable}

\begin{center}\rule{0.5\linewidth}{0.5pt}\end{center}

\section{6. Reproducibility}\label{reproducibility}

All handoff code, data, and bundled outputs sit under \texttt{handoff/}.
Setup:

\begin{Shaded}
\begin{Highlighting}[]
\BuiltInTok{cd}\NormalTok{ handoff}
\ExtensionTok{python} \AttributeTok{{-}m}\NormalTok{ venv .venv}
\BuiltInTok{source}\NormalTok{ .venv/bin/activate}
\ExtensionTok{pip}\NormalTok{ install }\AttributeTok{{-}r}\NormalTok{ requirements.txt}
\end{Highlighting}
\end{Shaded}

Verify the package:

\begin{Shaded}
\begin{Highlighting}[]
\ExtensionTok{pytest} \AttributeTok{{-}q}
\ExtensionTok{python}\NormalTok{ scripts/verify\_artifacts.py}
\ExtensionTok{python}\NormalTok{ scripts/bootstrap\_standalone.py}
\end{Highlighting}
\end{Shaded}

\textbf{Notebooks to open} (kernel cwd = \texttt{handoff/} or
\texttt{handoff/notebooks/}):

\begin{itemize}
\tightlist
\item
  \texttt{notebooks/GridDAM.ipynb}
\item
  \texttt{notebooks/VisionExperiments.ipynb}
\end{itemize}

\textbf{Bundled results} live in \texttt{results/} with checksums in
{\ttfamily\nolinkurl{results/MANIFEST.json}}. With \texttt{RUN\_EXPERIMENTS\ =\ False}
(vision) or \texttt{LOAD\_EXISTING\ =\ True} (grid multiscale),
notebooks reload those files without rerunning sweeps.

\textbf{PDF export} (Pandoc + memoir + Tectonic):

\begin{Shaded}
\begin{Highlighting}[]
\FunctionTok{bash}\NormalTok{ latex\_exports/compile\_memo.sh}
\end{Highlighting}
\end{Shaded}

Output: \texttt{latex\_exports/pdf/REPORT.pdf}.

\begin{center}\rule{0.5\linewidth}{0.5pt}\end{center}

\section{7. Scripts reference}\label{scripts-reference}

\subsection{7.1 Grid scripts}\label{grid-scripts}

\begin{longtable}[]{@{}
  >{\raggedright\arraybackslash}p{(\linewidth - 4\tabcolsep) * \real{0.32}}
  >{\raggedright\arraybackslash}p{(\linewidth - 4\tabcolsep) * \real{0.36}}
  >{\raggedright\arraybackslash}p{(\linewidth - 4\tabcolsep) * \real{0.32}}@{}}
\toprule\noalign{}
\begin{minipage}[b]{\linewidth}\raggedright
Script
\end{minipage} & \begin{minipage}[b]{\linewidth}\raggedright
Purpose \& key arguments
\end{minipage} & \begin{minipage}[b]{\linewidth}\raggedright
Default output
\end{minipage} \\
\midrule\noalign{}
\endhead
\bottomrule\noalign{}
\endlastfoot
{\ttfamily run\_\allowbreak grid\_\allowbreak multiscale.py} & Part 4 multiscale sweep.
\texttt{-\/-quick}, \texttt{-\/-scale-min}, \texttt{-\/-scale-max},
\texttt{-\/-n-values}, \texttt{-\/-k-values}, encoder/DAM hyperparams &
{\ttfamily\nolinkurl{results/grid/multiscale/}} \\
{\ttfamily run\_\allowbreak grid\_\allowbreak breakit.py} & Part 4b sensitivity.
\texttt{-\/-experiment} (orientations, base\_freq, modules, cells,
distributions, all), \texttt{-\/-quick} &
{\ttfamily\nolinkurl{results/grid/breakit/}} \\
\end{longtable}

Example quick smoke:

\begin{Shaded}
\begin{Highlighting}[]
\BuiltInTok{cd}\NormalTok{ handoff}
\BuiltInTok{source}\NormalTok{ .venv/bin/activate}
\ExtensionTok{python}\NormalTok{ scripts/run\_grid\_multiscale.py }\AttributeTok{{-}{-}quick}
\ExtensionTok{python}\NormalTok{ scripts/run\_grid\_breakit.py }\AttributeTok{{-}{-}experiment}\NormalTok{ orientations }\AttributeTok{{-}{-}quick}
\end{Highlighting}
\end{Shaded}

\subsection{7.2 Vision scripts}\label{vision-scripts}

\begin{longtable}[]{@{}
  >{\raggedright\arraybackslash}p{(\linewidth - 4\tabcolsep) * \real{0.32}}
  >{\raggedright\arraybackslash}p{(\linewidth - 4\tabcolsep) * \real{0.36}}
  >{\raggedright\arraybackslash}p{(\linewidth - 4\tabcolsep) * \real{0.32}}@{}}
\toprule\noalign{}
\begin{minipage}[b]{\linewidth}\raggedright
Script
\end{minipage} & \begin{minipage}[b]{\linewidth}\raggedright
Purpose \& key arguments
\end{minipage} & \begin{minipage}[b]{\linewidth}\raggedright
Default output
\end{minipage} \\
\midrule\noalign{}
\endhead
\bottomrule\noalign{}
\endlastfoot
{\ttfamily run\_\allowbreak vision\_\allowbreak synthetic.py} & Synthetic baseline.
\texttt{-\/-task-mode} (mixed, color\_only, position\_only),
\texttt{-\/-quick}, \texttt{-\/-device} &
{\ttfamily\nolinkurl{results/vision/synthetic/}} \\
{\ttfamily run\_\allowbreak vision\_\allowbreak model\_\allowbreak comparison.py} & Multi-backbone synthetic
sweep. \texttt{-\/-models} (or quick ViT-only), \texttt{-\/-device} &
{\ttfamily\nolinkurl{results/vision/model\_comparison/}} \\
{\ttfamily run\_\allowbreak vision\_\allowbreak naturalistic\_\allowbreak baseline.py} & Peterson baseline.
\texttt{-\/-category} (fruits, vegetables, animals),
\texttt{-\/-include-clip}, \texttt{-\/-quick} &
{\ttfamily\nolinkurl{results/vision/naturalistic/}} \\
{\ttfamily run\_\allowbreak vision\_\allowbreak naturalistic\_\allowbreak dam.py} & DAM on naturalistic
features. \texttt{-\/-category}, \texttt{-\/-baseline-dir},
\texttt{-\/-max-stage-a-configs}, \texttt{-\/-quick} &
{\ttfamily\nolinkurl{results/vision/naturalistic\_dam/}} \\
{\ttfamily run\_\allowbreak vision\_\allowbreak animals\_\allowbreak graphs.py} & \S{}5A graph regen.
\texttt{-\/-max-settings}, \texttt{-\/-max-seeds},
\texttt{-\/-max-layers}, \texttt{-\/-quick} &
{\ttfamily\nolinkurl{results/image\_transformers/}} \\
{\ttfamily run\_\allowbreak vision\_\allowbreak animals\_\allowbreak similarity\_\allowbreak ordered.py} & \S{}5B step 1. Same
caps as graphs script &
{\ttfamily\nolinkurl{results/image\_transformers\_similarity\_ordered/}} \\
{\ttfamily run\_\allowbreak vision\_\allowbreak animals\_\allowbreak similarity\_\allowbreak ordered\_\allowbreak fixed.py} & \S{}5B step
2 reorder plots. \texttt{-\/-input-dir}, \texttt{-\/-output-dir} &
{\ttfamily\nolinkurl{results/image\_transformers\_similarity\_ordered\_fixed/}} \\
\end{longtable}

Example vision smoke:

\begin{Shaded}
\begin{Highlighting}[]
\ExtensionTok{python}\NormalTok{ scripts/run\_vision\_synthetic.py }\AttributeTok{{-}{-}quick}
\ExtensionTok{python}\NormalTok{ scripts/run\_vision\_model\_comparison.py }\AttributeTok{{-}{-}models}\NormalTok{ vit\_base\_patch16\_224}
\ExtensionTok{python}\NormalTok{ scripts/run\_vision\_naturalistic\_baseline.py }\AttributeTok{{-}{-}category}\NormalTok{ fruits }\AttributeTok{{-}{-}quick}
\end{Highlighting}
\end{Shaded}

Full naturalistic DAM (slow):

\begin{Shaded}
\begin{Highlighting}[]
\ExtensionTok{python}\NormalTok{ scripts/run\_vision\_naturalistic\_baseline.py }\AttributeTok{{-}{-}category}\NormalTok{ animals }\AttributeTok{{-}{-}include{-}clip}
\ExtensionTok{python}\NormalTok{ scripts/run\_vision\_naturalistic\_dam.py }\AttributeTok{{-}{-}category}\NormalTok{ animals}
\end{Highlighting}
\end{Shaded}

\begin{center}\rule{0.5\linewidth}{0.5pt}\end{center}

\backmatter
\end{document}
